% !TEX program = pdflatex
\documentclass[11pt,a4paper]{article}

\usepackage[T1]{fontenc}
\usepackage[utf8]{inputenc}
\usepackage{lmodern}
\usepackage{microtype}

\usepackage{amsmath,amssymb,bm}
\usepackage{graphicx}
\usepackage{booktabs}
\usepackage{hyperref}

\title{Microstructure-Constrained Cloak Optimisation:\\
Pipeline and Validation Gap}
\author{(project notes)}
\date{\today}

\begin{document}
\maketitle

\section{Motivation}
Free per-cell $(\lambda,\mu,\rho)$ optimisation reaches near-perfect cloaking
($\langle|u|\rangle/\langle|u_{\rm ref}|\rangle \approx 1.0$) but produces
moduli that no real material realises. We constrain the optimisation to a
\emph{manifold of physically realisable microstructures} and then check
whether the corresponding microstructures, run at the pixel level, recover
the homogenised performance.

\section{Pipeline}

\paragraph{Step 1 -- CA dataset.}
$N{\sim}10^{6}$ binary $50\times 50$ unit cells generated by a
squared-assembly cellular automaton (cement / void).

\paragraph{Step 2 -- FEM homogenisation.}
Periodic FEM on each unique cell yields effective Lam\'e parameters
$(\lambda_{\rm eff},\mu_{\rm eff})$, density $\rho_{\rm eff}$, and the full
stiffness tensor $C_{\rm eff}$. Stored in \texttt{stiffness.h5}.

\paragraph{Step 3 -- GMM prior.}
A Gaussian mixture is fit to the standardised
$(\lambda_{\rm eff},\mu_{\rm eff},\rho_{\rm eff})$ point cloud. A flat-top
threshold $\tau$ at the $q$-th log-density quantile defines the realisable
manifold.

\paragraph{Step 4 -- Constrained optimisation.}
Per-cell parameters $(\lambda_c,\mu_c,\rho_c)$ are optimised to minimise the
top-surface displacement loss, with a penalty
$\max(0,\tau-\log p_{\rm GMM}(\lambda_c,\mu_c,\rho_c))$ pushing each cell
into the manifold.

\paragraph{Step 5 -- Snap to dataset.}
After optimisation each cloak cell is replaced by the dataset entry whose
standardised $(\lambda,\mu,\rho)$ is closest in Euclidean distance.

\paragraph{Step 6 -- Pixel-level validation.}
The matched microstructures tile the cloak region; the FEM is re-solved
with material assigned at the \emph{pixel} level (cement vs.\ void) on a
refined mesh ($\sim$1 element per micro-pixel side).

\section{Diagnostic experiment}
To separate the \emph{matching} error from the
\emph{pixel-vs.-homogenised} error we add an intermediate step: replace
each cloak cell's optimised $(\lambda,\mu,\rho)$ by its matched dataset
$(\lambda,\mu,\rho)$ but keep the homogenised FEM (a single $C$ and $\rho$
per cell). The script
\texttt{scripts/frequency\_sweep\_matched\_homogenized.py} implements this.

\section{Results}
Run on \texttt{output/cell20\_cement\_init/} (20$\times$15 macro grid,
100 cloak cells, dataset \texttt{ca\_bulk\_squared}) at $f^*=2.0$:

\begin{table}[h]
\centering
\begin{tabular}{lcc}
\toprule
Configuration & $\langle|u|\rangle/\langle|u_{\rm ref}|\rangle$ &
mesh refinement \\
\midrule
Optimised (homogenised, free $\lambda,\mu,\rho$) & $\approx 1.00$ & 3 \\
Matched (homogenised, dataset $\lambda,\mu,\rho$) & $0.9385$ & 3 \\
Matched (homogenised, dataset $\lambda,\mu,\rho$) & $0.9544$ & 20 \\
Validated (pixel-level, real microstructure)      & $\approx 0.70$ & 25 \\
\bottomrule
\end{tabular}
\caption{Cloaking ratio at $f^{*}=2.0$. Mesh refinement does not close the
matched-vs-validated gap.}
\end{table}

\noindent
Matching diagnostics on the 100 cloak cells: 28 unique dataset entries
used; standardised L$^2$ distance median $0.012$, max $0.031$.
Per-cell relative shifts in physical units:
$|\Delta\lambda|/\lambda \approx 0.012$,
$|\Delta\mu|/\mu \approx 0.899$,
$|\Delta\rho|/\rho \approx 0.001$.
The standardisation hides a near-100\% mismatch in $\mu$ because the
dataset's spread in $\mu$ is small relative to that of $\lambda$ and
$\rho$.

\section{Discussion}

\paragraph{The $\sim$5\% matching gap is benign.}
Going from the free optimum (1.00) to the matched homogenised material
(0.94--0.95) costs about five percentage points. This is the price of
projecting onto the discrete dataset manifold, and it is straightforwardly
improved by (i) enlarging the dataset, (ii) reweighting the matching
distance so $\mu$ is not nearly ignored, or (iii) replacing the nearest-
neighbour snap with a soft assignment / generative model.

\paragraph{The $\sim$25\% pixel-vs-homogenised gap is the real problem.}
The matched homogenised cloak retains $\sim$95\% of the optimum, but
once the same dataset entries are inserted as actual microstructures the
performance collapses to $\sim$70\%. Mesh refinement does not close this
gap, so it is physical, not numerical. Candidate causes (to investigate):

\begin{enumerate}
  \item \textbf{Hidden anisotropy.} The descriptor $(\lambda,\mu,\rho)$ is
        an isotropic projection of an effective stiffness tensor that is
        in general anisotropic (up to four independent moduli in 2D
        orthotropy). Two microstructures with identical $(\lambda,\mu)$
        but different orientation of the principal axes scatter waves
        very differently. The matched homogenised solver sees only the
        isotropic projection; the pixel solver sees the true anisotropic
        medium.
  \item \textbf{Static vs.\ dynamic homogenisation.} The dataset is built
        from quasi-static periodic FEM. Scale separation
        $\lambda_{\rm wave}\gg\ell_{\rm micro}$ is only marginally
        satisfied at $f^{*}=2$ (the wavelength in cement is $\sim$1
        macro cell, and intra-cell features can be a few pixels).
        Dynamic effective properties (Willis-type couplings,
        frequency-dependent moduli) are then non-negligible and absent
        from the lookup.
  \item \textbf{Scattering at void interfaces.} With void ratio $10^{-6}$
        every cement/void boundary acts as a free surface. The pixel
        cloak contains $O(10^{4})$ such interfaces, each a stress
        concentrator and partial reflector that no bulk descriptor can
        encode.
  \item \textbf{Inter-cell discontinuities.} Adjacent macro cells in the
        homogenised model carry smoothly varying $(\lambda,\mu,\rho)$;
        in the pixel model their microstructures abut without
        topological continuity, producing extra scattering at the
        cell-grid lines.
\end{enumerate}

\paragraph{Proposed disambiguating experiment.}
For each cloak cell, run a fresh static periodic FEM on its actual matched
microstructure and fit the \emph{full} effective $C$ tensor (anisotropic).
Insert these per-cell $C$ tensors in the homogenised solver. If the result
recovers $\sim$0.95, the gap is essentially the isotropic-truncation of
the descriptor (Hypothesis 1) and can be repaired by storing and
optimising over anisotropic moduli. If it still drops to $\sim$0.70, the
remainder is genuine finite-wavelength / scattering physics
(Hypotheses 2--4) and can only be addressed by either dynamic
homogenisation or by optimising directly in microstructure space.

\end{document}
